\newpage
\section{Conclusion}

Ce stage avait pour objectif d'évaluer la qualité de contextes causaux contrefactuels dans le cadre des Systèmes de Transitions d'États Étiquetés, en tenant compte des normes. Une réflexion a notamment été menée sur la façon de représenter ces normes au sein de ce formalisme : le choix s'est porté sur une pondération des transitions du LTS. Ce travail répond à cet objectif au travers de deux contributions principales.

La première est une mesure d'évaluation constituée d'une fonction d'attribution de coût modulaire couplée à des filtres a posteriori définissant une forme d'acceptabilité des contrefactuels. La fonction de coût combine une mesure de distance à l'observation basée sur la distance de Levenshtein enrichie par la notion de mutabilité et une mesure de vraisemblance normative. Cette vraisemblance peut être calculée par comparaison à des politiques optimales ou par agrégation des pondérations le long d'une trace. Ce type de mesure, intégrant des normes sur des objets aussi riche que les LTS, qui distinguent explicitement états et événements et permettent de représenter des chaînes causales complexes, n'existait pas dans la littérature. La modularité de la fonction proposée est un atout important et permet d'adapter finement son usage en fonction de l'application : choix des opérateurs d'agrégation, des pondérations ou filtres a posteriori.

L'application à l'exemple du poison développée dans la section \ref{ex_poison} illustre la pertinence de l'approche proposée : les différentes configurations testées produisent des classements de contrefactuels cohérents avec l'intuition attendue de chacun des critères. Cette approche a également été testée sur d'autres exemples, tels que l'exemple du jardinier, confirmant les résultats obtenus.

La seconde contribution est un parseur permettant de construire automatiquement un LTS à partir d'un problème PDDL, dont a été démontré formellement la correction, la terminaison ainsi que la représentabilité. Ce parseur nous a permis de générer des exemples classiques de la littérature pour valider notre approche.

Ce travail a nécessité d'articuler des concepts issus de plusieurs domaines de recherche tels que la causalité effective, l'explicabilité en IA et les sciences cognitives dans un même cadre formel. Cette démarche a mené à l'identification de divergences sémantiques se trouvant derrière des définitions propres à chaque domaine, comme, la notion d'acceptabilité. Elle a aussi permis d'intégrer la notion de mutabilité, issue des sciences cognitives à la définition de la distance à l'observation. 

Le travail formel est accompagné d'une implémentation complète en Python, comprenant le cadre théorique (LTS, contextes causaux, fonction de coût modulaire et ses différentes composantes) ainsi que le processus d'évaluation incluant le parseur de PDDL vers LTS. Cette implémentation a permis de tester les différentes configurations de la fonction de coût sur des exemples, et constitue une base réutilisable pour de futurs travaux s'appuyant sur ce cadre.


Ce travail ouvre plusieurs perspectives. Sur le plan formel, le cadre pourrait être étendu à des transitions stochastiques, ce qui permettrait de mieux représenter la réalité, mais soulèverait en contrepartie des questions de stabilité et de robustesse des politiques face à l'incertitude, nécessitant l'ajout de nouveaux termes dans la fonction de coût pour en tenir compte. Une seconde piste serait d'intégrer un terme causal explicite, prenant en compte les fluents décrivant les états. On pourrait également envisager une évaluation floue des normes, plutôt qu'une pondération numérique stricte afin de mieux représenter le respect ou la violation de norme. Enfin, maintenant que nous avons défini une fonction de coût permettant d'évaluer la qualité de contextes causaux contrefactuels, il serait pertinent de nous intéresser à la construction automatique de contrefactuels visant à en minimiser la valeur.



% -synthèse résultats du poison ce que j'ai fait permet de bien traiter cet exemple et corresponds à ce qu'on attent intuitivement pour chaque des choix

% -j'ai defini une mesure qui prends en compte des normes sur des objets complexes que sont les LTS, ce type de mesure n'existait pas

% - j'ai introduit un parseur ou je verifie formellement ses propriétés et il me permet de generer des exemples

% - je peux/j'ai appliquer mon programme sur les exemples classiques de causalité 

% -la fonction est modulaire et permet une grande liberté on peut plug ce qu'on veut en fonction de l'applications

% -reprendre les objectifs de l'Introduction

% -j'ai fais des liens entre des domaines et je les ai integré à mon travail

% -rajouter les trucs remarquables/interessant/nouveau que j'ai fait

% -perspectives :

% Intégration de transitions stochastiques dans les LTS : critère de stabilité et robustesse

%              Terme causal : prise en compte des fluents

%             Évaluation floue des normes

%              Construction automatique de politiques

        